\documentclass[a4paper,12pt]{extarticle}
\usepackage[T2A]{fontenc}
\usepackage[utf8]{inputenc}
\usepackage[english,russian]{babel}
\usepackage[rus]{datetime2}          % русская локализация
\usepackage{etoolbox}                % для удобных манипуляций
\usepackage{amsmath, amssymb}
\usepackage{enumitem}
\usepackage{tikz}
\usepackage{graphicx}
\usepackage{stackengine}

\usepackage{fancyhdr}
\pagestyle{fancy}
\fancyhf{} % очистить все поля
\fancyhead[L]{\leftmark}   % слева вверху – текущий раздел (section)
\fancyhead[R]{\thepage}    % справа вверху – номер страницы
\renewcommand{\headrulewidth}{0.4pt} % тонкая линия под заголовком
\setlength{\headwidth}{\paperwidth}
\fancyheadoffset{0pt}   % или \fancyhfoffset{0pt}

\usepackage{hyperref}
\hypersetup{
    linktocpage,
    colorlinks=true
}

\usepackage{makecell}
\topmargin=-2cm
\oddsidemargin=-2cm
\textheight=27.5cm
\textwidth=19cm
\parindent=0em
\parskip=0mm


% Базовые символы(теорема, лемма и др.)
\newcommand{\TheoMark}{%
	\hspace{5.5pt}%
  \scalebox{2.3}{\stackengine{3.85pt}{\scalebox{1.5}{$\triangledown$}}{\scriptsize \textnormal{T}}{O}{c}{F}{T}{L}}%
}
\newcommand{\lemmaMark}{%
  	\hspace{5.4pt}%
    \scalebox{1.5}{\stackinset{c}{0pt}{c}{0pt}{\scalebox{1.43}{$\Lambda$}}{\scalebox{1.5}{$\bigcirc$}}}%
}
\newcommand{\DefMark}{%
    \scalebox{2}{\underline{O}}%
}


% Окружения для часто появляющихся конструкций
\newenvironment{theorem}[1][]{%
  \par\medskip
  \noindent
  \begin{tabular}{@{}p{0.03\textwidth} p{0.97\textwidth}@{}}
    \raggedright \makebox[0.03\textwidth][c]{\TheoMark} & 
    \begin{minipage}[t]{\linewidth}
      \textbf{%
        \if\relax\detokenize{#1}\relax\else #1\fi
      }%
      \par\nobreak\noindent\ignorespaces
}{%
      \end{minipage}
  \end{tabular}%
  \par
}

\newenvironment{lemma}[1][]{%
  \par\medskip
  \noindent
  \begin{tabular}{@{}p{0.03\textwidth} p{0.97\textwidth}@{}}
    \raggedright \makebox[0.03\textwidth][c]{\lemmaMark} & 
    \begin{minipage}[t]{\linewidth}
      \textbf{%
        \if\relax\detokenize{#1}\relax\else #1\fi
      }%
      \par\nobreak\noindent\ignorespaces
}{%
      \end{minipage}
  \end{tabular}%
  \par\medskip
}

\newenvironment{definition}[1][]{%
  \par\medskip
  \noindent
  \begin{tabular}{@{}p{0.03\textwidth} p{0.97\textwidth}@{}}
    \raggedright \DefMark & 
    \begin{minipage}[t]{\linewidth}
      \textbf{%
        \if\relax\detokenize{#1}\relax\else #1\fi
      }%
      \par\nobreak\noindent\ignorespaces
}{%
      \end{minipage}
  \end{tabular}%
  \par\medskip
}

\newenvironment{proof}{%
  \par\medskip
  \noindent
  $\blacktriangleright$\
  \par\nobreak\noindent\ignorespaces
}{%
  \par\nobreak
  \noindent
  \hfill$\blacksquare$
  \par\medskip
}

\newenvironment{note}[1][]{\par\medskip\textbf{Замечание\relax #1. }}{\par\medskip}

\usepackage{iftex}

\ifPDFTeX
  \newcommand{\compilerinfo}{\pdftexbanner}
\else\ifLuaTeX
  \newcommand{\compilerinfo}{\luatexbanner}
\else\ifXeTeX
  \newcommand{\compilerinfo}{XeTeX, версия \the\XeTeXversion\ (rev. \the\XeTeXrevision)}
\else
  \newcommand{\compilerinfo}{(неизвестный движок)}
\fi\fi\fi

\DTMnewdatestyle{russiandate}{%
  \renewcommand{\DTMdisplaydate}[4]{%
    \number##3\space \DTMmonthname{##2}\space ##1\space г.%
  }%
}
\DTMsetdatestyle{russiandate}

\DTMlangsetup{showdayofmonth=false}

\newcommand{\currenttime}{%
  \DTMcurrenttime%
}

% Переопределяем \date
\date{Скомпилирован: \today\ в \currenttime \\ \compilerinfo}

\begin{document}
\title{Уравнения математической физики\\Лекционный материал}
\maketitle

\tableofcontents
\break

\section{Лекция 1: Уравнения в частных производных 2-го порядка}
Многие задачи механики сплошной среды и физики приводят к исследованию дифференциальных уравнений с частными производными второго порядка (ДУЧП2П). Так, например.

1) При изучении различных видов волн (упругих, звуковых, электромагнитных) мы приходим к \textsl{волновому уравнению}:
\[
\frac{\partial^2 u}{\partial t^2} = c^2 \left( \frac{\partial^2 u}{\partial x^2} + \frac{\partial^2 u}{\partial y^2} + \frac{\partial^2 u}{\partial z^2} \right) + f(x, y, z, t),
\]
где $c$ -- скорость распространения волн в данной среде, $f$ -- источник.

2) Процессы распространения тепла в однородном изотропном теле описываются уравнением теплопроводности
\[
\frac{\partial u}{\partial t} = a^2 \left( \frac{\partial^2 u}{\partial x^2} + \frac{\partial^2 u}{\partial y^2} + \frac{\partial^2 u}{\partial z^2} \right) + f(x, y, z, t),
\]
где $a^2$ – коэффициент температуропроводности данной среды.

3) При рассмотрении установившегося теплового состояния в однородном изотропном теле мы приходим к уравнению Пуассона:
\[
	\frac{\partial^2 u}{\partial x^2} + \frac{\partial^2 u}{\partial y^2} + \frac{\partial^2 u}{\partial z^2} = -\tilde{f}(x, y, z),
\]
где $\tilde{f}(x, y, z) = \frac{f(x, y, z)}{a^2}$

При отсутствии источников тепла внутри тела уравнение теплопроводности переходит в уравнение Лапласа:
\[
	\frac{\partial^2 u}{\partial x^2} + \frac{\partial^2 u}{\partial y^2} + \frac{\partial^2 u}{\partial z^2} = 0,
\]

Волновое уравнение (ВУ), уравнение теплопроводности (УТ), уравнение Пуассона (УП) и уравнение Лапласа (УЛ) часто называют основными уравнениями математической физики.

\bigskip

Каждое из основных уравнений математической физики имеет бесчисленное множество частных решений. При решении конкретной физической задачи, необходимо из всех этих решений выбрать то, которое удовлетворяет некоторым дополнительным условиям, вытекающим из её физического смысла. Итак, задачи математической физики состоят в отыскании решений уравнений в частных производных, удовлетворяющим некоторым дополнительным условиям. Такими дополнительными условиями чаще всего являются так называемые граничные (краевые) условия, т.е. условия, заданные на границе рассматриваемой среды, и начальные условия, относящегося к одному какому-нибудь моменту времени, с которого начинается изучение данного физического явления.

\bigskip

Обратим внимание на одно важное обстоятельство.

\textsl{Задача математической физики считается поставленной корректно (в смысле Адамара), если решение задачи, удовлетворяющее всем её условиям, существует, единственно и устойчиво} (непрерывно зависит от данных (параметров) задачи, т.е. малые изменения любого из данных задачи должно вызывать соответственно малое изменение решения). Требование устойчивости выражает физическую определенность поставленной задачи.

\bigskip

Дадим необходимые определения.

\begin{definition}[Дифференциальное уравнение в частных производных второго порядка (ДУЧП2П)]
Пусть $G$ --- область в $\mathbb{R}^n$, x=($x_1$, \ldots, $x_n$) --- точка из $G$. Дифференциальным уравнением в частных производных второго порядка (ДУЧП2П) с независимыми переменными $x_1, \ldots, x_n$ называется соотношение, связывающее значение неизвестной функции n переменных $u=u(x_1, \ldots, x_n)$, и её частных производных 1\textsuperscript{го} и 2\textsuperscript{го} порядков:
\[
\frac{\partial u}{\partial x_1}, \ldots, \frac{\partial u}{\partial x_n},
\frac{\partial^2 u}{\partial^2 x_1}, \ldots, \frac{\partial^2 u}{\partial x_i \partial x_j}, \ldots \frac{\partial^2 u}{\partial^2 x_n}
\] вида:
\[
F(x_1, \ldots, x_n, u, \frac{\partial u}{\partial x_1}, \ldots, \frac{\partial u}{\partial x_n},
\frac{\partial^2 u}{\partial^2 x_1}, \ldots, \frac{\partial^2 u}{\partial x_i \partial x_j}, \ldots \frac{\partial^2 u}{\partial^2 x_n})=0
\]
\end{definition}

Решением ДУЧП2П называется функция  класса , которая удовлетворяет данному уравнению.

Уравнение называется \textsl{линейным} относительно старших производных – производных второго порядка (ДУЧП2ПЛОСП), если оно имеет вид:

\[
\sum_{i,j=1}^{n} a_{ij}(x)\frac{\partial^2u}{\partial x_i \partial x_j} + \Phi(x, u, \operatorname{grad}(u)) = 0
\]
где коэффициенты при старших (вторых) производных $a_{ij}(x) \in C(G) (i, j = 1, \ldots,n)$ являются функциями только независимых переменных $x_1, \ldots, x_n$; 
а $\operatorname{grad}(u) = (\frac{\partial u}{\partial x_1}, \ldots, \frac{\partial u}{\partial x_n})$ --- градиент функции $u(x) = u(x_1, \ldots, x_n)$.

\bigskip

Коэффициенты при старших (вторых) производных образуют квадратную матрицу $A(x) = (a_{ij}(x))$  – матрицу старших коэффициентов:



\begin{note}
$\operatorname{grad}(u) = (\frac{\partial u}{\partial x_1}, \ldots, \frac{\partial u}{\partial x_n})$.
\end{note}
\begin{note}
Если $\Phi$ - линейная относительно $u(x)$, $\operatorname{grad}(u)$, то получаем линейное уравнение в частных производных 2\textsuperscript{го} порялка(ЛУЧП2П):
$$\sum_{i,j=1}^{n} a_{ij}(x)\frac{\partial^2u}{\partial x_i \partial x_j} + \sum_{i=1}^{n} b_{i}(x)\frac{\partial u}{\partial x_i} + c(x)u = f(x),$$
\[
  x \in G,\quad a_{ij}(x), b_i(x), c(x) \in C(G); \qquad\qquad (i,j = \overline{1,n})
\]
$$f(x) \in C(G)$$
\end{note}
Иногда пишут:
$L[u]=f(x)$, где
$$\displaystyle L= \sum_{i,j=1}^{n} a_{ij}(x)\frac{\partial^2}{\partial x_i \partial x_j} + \sum_{i=1}^{n} b_{i}(x)\frac{\partial }{\partial x_i} + c(x)$$
$\displaystyle L: C^2(G)\mapsto C(G)$ -- линейный дифференциальный оператор 2\textsuperscript{го} порядка(ЛДО2П).\par\medskip


\label{abbr:LDUChPSPPK}
\newcommand{\LDUChPSPPK}{\hyperref[abbr:LDUChPSPPK]{ЛДУЧП2ППК}}
Если же все коэффициенты линейного дифференциального уравнения частных производных 2\textsuperscript{го} порядка(ЛДУЧП2П): не зависят от $\vec x$, то получаем ЛДУЧП2П с постоянными коэффициентами (\LDUChPSPPK{}):

\[
\sum_{i,j=1}^{n} a_{ij}\frac{\partial^2u}{\partial x_i \partial x_j} + \sum_{i=1}^{n} b_{i}\frac{\partial u}{\partial x_i} + cu = f(x), \quad x \in G,
\]
где \(a_{ij}, b_i, c\) --- коэффициенты, \(i,j = \overline{1,n}\), \(f(x) \in C(G)\) --- правая часть. Это уравнение можно записать как \(L[u]=f\).\\
Здесь
$\displaystyle L= 
  \sum_{i,j=1}^{n} a_{ij}\frac{\partial}{\partial x_i} \frac{\partial}{\partial x_j}
  +\sum_{i=1}^{n} b_{i}\frac{\partial }{\partial x_i} + c$.\\ Первая сумма это \emph{квадратичная форма}, вторая -- \emph{линейная форма}.
\subsection{Линейная невырожденная замена переменных}
Осуществим в \LDUChPSPPK{} линейную невырожденную замену переменных:\\
$\displaystyle\vec x=\vec x(\vec \xi),$ где $\displaystyle\vec \xi=(\xi_1, \ldots, \xi_n),\quad \xi_k = \sum_{i=1}^{n}s_{ik}x_i,\ (k=\overline{1,n})$.\\
В матричной форме: $$\vec \xi = S^T \vec x$$
S --- матрица коэффициентов: S = ($s_{ik}$)$\Rightarrow$

\[
\begin{cases}
\xi_1 = s_{11}x_1 + \ldots + s_{n1}x_n \\
\ldots \\
\xi_i = s_{1i}x_1 + \ldots + s_{ni}x_n \\
\ldots \\
\xi_n = s_{1n}x_1 + \ldots + s_{nn}x_n
\end{cases}\Longrightarrow
\quad \det S \neq 0 \Longrightarrow \vec x=(S^{-1})^T\vec\xi
\]
Обозначим
\[
u(x(\xi))=\tilde{u}(\xi) \Rightarrow \tilde{u}(\xi)=u(x)\]
\[
\text{Т.к. } \frac{\partial \xi_k}{\partial x_i} = s_{ik} \text{, то} \frac{\partial u}{\partial x_i} = \sum_{k=1}^n\frac{\partial \tilde{u}}{\partial \xi_i}\frac{\partial \xi_k}{\partial x_i} = \sum_{k=1}^ns_{ik}\frac{\partial \tilde{u}}{\partial \xi_k}
\]


\begin{note}
Из $\displaystyle\frac{\partial }{\partial x_i} = \sum_{k=1}^n s_{ik}\frac{\partial}{\partial \xi_k}$ имеем:
\[\frac{\partial^2u}{\partial x_i \partial x_j}=\frac{\partial}{\partial x_i}(\frac{\partial u}{\partial x_j})=\sum_{l,k=1}^{n}\frac{\partial^2 \tilde{u}}{\partial\xi_k \partial\xi_l}s_{ik}s_{jl}\]
Подставляя полученные выражения для $\frac{\partial u}{\partial x_i}$ и $\frac{\partial^2 u}{\partial x_i \partial x_j}$ в \LDUChPSPPK{}, получим:

\[\xi_k=\sum_{i=1}^{n}s_{ik}x_i,\ \xi=S^Tx, x=(S^{-1})^T\xi\Rightarrow\]

\[
\sum_{k,l=1}^n \frac{\partial^2 \tilde{u}}{\partial \xi_k \partial \xi_l} \sum_{i,j=1}^n a_{ij}s_{ik}s_{jl} + \sum_{k=1}^n \frac{\partial \tilde{u}}{\partial \xi_k}\sum_{i=1}^nb_is_{ik}+c\tilde{u}=\tilde{f}
\]
Обозначим через 
\[\tilde{a}_{kl}=\sum_{i,j=1}^na_{ij}s_{ik}s_{jl}, \quad\tilde{b}_k=\sum_{i=1}^nb_is_{ik}\]
получаем:
$\displaystyle
\tilde{L}\tilde{u}=\sum_{k,l=1}^n\tilde{a}_{kl}\frac{\partial^2\tilde{u}}{\partial\xi_k\partial\xi_l}+\sum_{k=1}^n\tilde{b}_k\frac{\partial\tilde{u}}{\partial\xi_k}+c\tilde{u}=\tilde{f}(\xi),
$\\
где $\tilde{f}(\xi)=f(x(\xi))$\\
При этом матрица $A=(a_{ik})$ старших коэффициентов исходного ЛУЧП2ППК преобразуется по формулам\\
$
\displaystyle
  \tilde{A}=S^TAS 
$
   (из линейной алгебры)\\
Т.е. при невырожденной линейной замене переменных:
\[
\xi=S^Tx, detS\ne0,
\] матрица A преобразуется по формулам
\[
K=\sum_{i,j=1}^na_{ij}y_iy_j=y^TAy,\ y=(y_1, \ldots, y_n)^T
\]
\end{note}

\begin{note}
считаем A --- симметричной ($A^T=A$)
\end{note}


\section{Лекция 2: Линейная невыр. замена независимых переменных}
Рассмотрим новый дифференциальный оператор:\\
$$\tilde{L}=\sum_{k,l=1}^{n} \tilde{a}_{kl}\frac{\partial}{\partial \xi_k} \frac{\partial}{\partial \xi_l}
+\sum_{k=1}^{n} \tilde{b}_{k}\frac{\partial }{\partial \xi_k} + \tilde{c}$$

Обозначим: $A = (a_{ij})$ и $\tilde{A} = (\tilde{a}_{ij})$ -- матрицы старших коэффициэнтов(МСК), где $A^T=A$ и $\tilde{A}^T=\tilde{A}$

В силу того, что $\displaystyle u(x)\in C^2(G)$, выполняется $\displaystyle \frac{\partial^2 u}{\partial x_i \partial x_j}= \frac{\partial^2 u}{\partial x_j \partial x_i}$. А значит:
$$\sum_{i,j=1}^{n} a_{ij}\frac{\partial^2 u}{\partial x_i\partial x_j}=
  \sum_{i,j=1}^{n} a_{ji}\frac{\partial^2 u}{\partial x_i\partial x_j}=
  \sum_{i,j=1}^{n} \frac{a_{ji}+a_{ij}}{2}\cdot\frac{\partial^2 u}{\partial x_i\partial x_j}$$

Таким образом, для ЛДУЧП2ППК можно заменить МСК $A$ на симметричную, вида: $\displaystyle A_c=\frac{A+A^T}{2}$ \\
Тогда при ЛНЗНП($\displaystyle\xi=S^T x$) МСК меняется следующим образом: $\displaystyle\tilde{A}=S^T A S$\\
Сопоставим для МСК $A$ квадратичную форму $K=\sum_{i,j=1}^n a_{ij}y_i y_j$\\
Квадратичную форму $K$ можно привести к каноническому виду:\\

\begin{tabular}{p{0.35\textwidth} p{0.6\textwidth}}
  \hspace{1cm}\(\displaystyle 
  K=\sum_{k=1}^{p}\eta^2_k-\sum_{k=p+1}^{r=p+q}\eta^2_k,
  \) &
  $p$ и $q$ -- положительные и отрецательные индексы инерции для $K$; $r$ -- ранг; $\sigma=p-q$ -- сигнатура $K$
\end{tabular}

\textbf{4 инварианта} $K\Longrightarrow$ при линейной невырожденной замене переменных $\xi = S^T x$ (т.е. $ y=S\eta $) приводит $K$  к каноническому виду $\displaystyle K = \eta^T \tilde{A} \eta$, тогда:
$$ 
  \tilde{L} \tilde{u} =
    \sum^p_{k=1}\frac{\partial^2\tilde u}{\partial \xi^2_k} -
    \sum^{p+q}_{k=p+1}\frac{\partial^2\tilde u}{\partial \xi^2_k} +
    \sum^{n}_{k=1}\frac{\partial\tilde u}{\partial \xi_k} + c \tilde u = \tilde f(\xi)
$$

\subsection{Классификация ЛДУЧП2ППК}

\textbf{I. Эллиптический.}
Если для квадратичной формы $K=y^T A y$, $r=n$ и в её каноническом виде все слагаемые одинакового знака($p=r=n,\,q=0$ либо $p=0,\,q=r=n$), то ЛДУЧП2ППК называется уравнением эллиптического типа.\smallskip\\
\emph{Пример:} Уравнение Лапласса: $\displaystyle \Delta u = 0$\smallskip

\textbf{II. Гиперболический.}
Если для квадратичной формы $K = y^TAy$, $r=n$ и в её каноническом виде $1\le p \le n-1,\,q \in [1,n-1]$, то ЛДУЧП2ППК называется уравнением гиперболического типа.\medskip\\
\emph{Пример:} Волновое уравнение: $\displaystyle \frac{\partial^2 u}{\partial t^2} = c^2 \Delta u + f(x, y, z, t)$
\begin{note}
  Если $p=1,\,q=n-1$ или $p=n-1,q=1$ -- нормально гиперболичкеский вид
\end{note}

\textbf{III. Параболический.}
Если для квадратичной формы $K = y^TAy$, $r<n$, то ЛДУЧП2ППК называется уравнением гиперболического типа.\medskip\\
\emph{Пример:} Уравнение теплопроводности: $\displaystyle \frac{\partial u}{\partial t} = c^2 \Delta u + f(x, y, z, t)$
\begin{note}
  Если $p=r,\,q=0$ или $p=0,q=r$ -- нормально параболический вид
\end{note}

Вернемся к ДУЧП2ПЛОСП:\\
\begin{tabular}{p{0.55\textwidth} p{0.35\textwidth}}
  \hspace{1cm}\(\displaystyle 
  \sum_{i,j=1}^{n} a_{ij}(x)\frac{\partial^2u}{\partial x_i \partial x_j} + \Phi(x_1, \ldots, x_n, u, \operatorname{grad}(u)) = 0,
  \) &
  $x \in G \subset \mathbb R^n$, $u \in C^2(G)$, $a_{ij} \in C(G)$
\end{tabular}\par\medskip

ЛДУЧП2П:%FIXME Че тут вообще произошло? Нить повествования потеряна
$$
Lu=\sum_{i,j=1}^{n} a_{ij}(x)\frac{\partial^2u}{\partial x_i \partial x_j} + \sum_{i=1}^{n} b_{i}(x)\frac{\partial u}{\partial x_i} + cu = f(x), \quad x \in G,
$$

Рассмотрим $\forall$ невырожденную замену переменных $\displaystyle \xi=\xi(x)\in C^2(G) : \frac{\partial(\xi_1,\dots,\xi_n)}{\partial(x_1,\dots,x_n)}\ne0 \Longrightarrow$ локально $\exists$ обратная замена $x=x(\xi)$ \\
Обозначим $\tilde u(\xi) = u(x(\xi)) \Longrightarrow u(x)=\tilde u(\xi(x))$. Тогда:
$$ 
  \frac{\partial u}{\partial x_i} = \sum^n_{k=1} \frac{\partial \tilde u}{\partial \xi_k}\cdot\frac{\partial \xi_k}{\partial x_i} 
  \Longrightarrow \frac{\partial^2 u}{\partial x_i \partial x_j}=
      \sum^n_{k=1} \frac{\partial^2 \tilde u}{\partial \xi_k\partial \xi_l}
      \cdot\frac{\partial \xi_k}{\partial x_i}\cdot\frac{\partial \xi_l}{\partial x_j}+
      \sum^n_{k=1} \frac{\partial \tilde u}{\partial \xi_k}\cdot\frac{\partial^2 \xi_k}{\partial x_i\partial x_j}
$$
или можно записать:
$$ 
  \frac{\partial}{\partial x_i} = \sum^n_{k=1} \frac{\partial \xi_k}{\partial x_i}\cdot\frac{\partial}{\partial \xi_k}
  \Longrightarrow \frac{\partial^2}{\partial x_i \partial x_j}=
      \sum^n_{k=1} \frac{\partial \xi_k}{\partial x_i}\cdot\frac{\partial \xi_l}{\partial x_j}
      \cdot\frac{\partial^2 }{\partial \xi_k\partial \xi_l}+
      \sum^n_{k=1} \frac{\partial^2 \xi_k}{\partial x_i\partial x_j}\cdot\frac{\partial}{\partial \xi_k}
$$
--- линейные операторы и нелинейные операторы




%FIXME


Подставляя выражение для $\frac{\partial u}{\partial x_i}$ и $\frac{\partial^2 u}{\partial x_i \partial x_j}$ в ЛДУЧП2ПЛОСП, получим:
$$ 
\sum_{k,l=1}^n \frac{\partial^2 \tilde u}{\partial \xi_k \partial \xi_l}( \sum_{i,j=1}^n a_{ij} \frac{\partial \xi_k}{\partial x_i}\cdot \frac{\partial \xi_l}{\partial x_j} )
+ \sum_{k=1}^n \frac{\partial \tilde u}{\partial\xi_k}(\sum_{i,j=1}^n a_{ij}\frac{\partial^2 \xi_k}{\partial x_i \partial x_j} + \sum^n_{i=1}b_i \frac{\partial \xi_k}{\partial x_i}) + c\tilde u = f
$$%FIXME big ()
Обозначим:
$$ \tilde a_{kl}(\xi) = \sum^n_{i,j = 1}a_{ij}\frac{\partial \xi_k }{\partial x_i}\cdot\frac{\partial \xi_l}{\partial x_j}; \quad k,l = \overline{1,n} $$

Тогда ДУЧП2ПЛОСП принимает вид:
$$ \sum^n_{k,l=1}\tilde a_{kl}(\xi) \frac{\partial^2 \tilde u}{\partial \xi_k \partial \xi_l} + \tilde \Phi(\xi, \tilde u, \operatorname{grad}(\tilde u)) = 0 $$

При этом ЛДУЧП2П принимает вид: 
$$ \tilde L \tilde u = \sum^n_{k,l=1}\tilde a_{kl} \frac{\partial^2\tilde u}{\partial \xi_k \partial \xi_l} +
\sum^n_{k=1} \tilde b_k (\xi) \frac{\partial \xi_k}{\partial x_i}, \qquad k=\overline{1,n}
$$
где: $$\displaystyle \tilde b_n(\xi) = \sum^n_{i,j=1}a_{i,j}(x)\frac{\partial^2 \xi_k}{\partial x_i\partial x_j}+\sum^n_{i=1}b_i(x)\frac{\partial\xi_k}{\partial x_i},\qquad k=\overline{1,n}$$\\
$$\tilde c(\xi)=c(x(\xi)),\qquad\tilde f(\xi) = f(x(\xi))$$

Считаем матрицу старших коэффициентов $A(x)=(a_{ij}(x))$ -- симметрической. Аналогично, $ \tilde A(\xi)=(\tilde a_{ij}) $ в ДУЧП2ПЛОСП(ЛДУЧП2П).\\
Фиксируем некоторую точку $x=x^{(0)}\in G$.\\
Обозначим $\displaystyle\xi^{(0)} = \xi(x^{(0)}) \Longrightarrow \tilde a_{kl}(\xi^{(0)}) = \sum^n_{i,j=1}a_{ij}(x^{(0)})\xi_{ik}\xi_{jl}$,
где $\displaystyle \xi_{ik} = \frac{\partial \xi_k}{\partial x_i}$ (и т.п.), или
$$ 
\tilde A(\xi^{(0)}) = S^T A(x^{(0)}) \Longrightarrow \sum^n_{k,l=1}\tilde a_{kl}\frac{\partial^2 \tilde u}{\partial \xi_k \partial \xi_l}+
\tilde\Phi(\xi,\tilde u,\operatorname{grad}(\tilde u))=0
$$
и в фиксированной точке $ x^{(0)} \in G \rightarrow A(\xi^{(0)}) = S^TA(x^{(0)})S $





















\section{Лекция 3: Классификация ДУЧП2ПЛОСП в точке}
Рассмотрим квадратичную форму вида:
$$ K=\sum_{i,j=1}^n a_{ij}(x^{(0)})y_i y_j = y^T A(x^{(0)})y, \qquad y = (y_1, \dots , y_n) $$
при ЛНЗНП:
\[
y=S\eta,\ \eta=(\eta_1, \ldots, \eta_n)^T,\ detS \neq 0:
\quad
K=y^TA(x^{(0)})y= \eta^T S^T A(x^{(0)}) S\eta =\eta^T\tilde A(\xi^{(0)})\eta
\]

Линейной невырожденной замене независимых переменных

\[y=S\eta\quad (y_i=\sum_{k=1}^ns_{ik}\eta_k)\] соответствует: \[\xi=S^T(x-x^{(0)});\qquad (\xi_k=\sum_{k=1}^ns_{ik}(x_i-x_i^{(0)})) \]

Распишем это подробнее:

Для любой квадратичной формы $K=y^TA(x^{(0)})y$ существует ЛНЗНП: $y=S\eta$, приводящая эту форму к каноническому виду:
\[
  K = \sum_{k=1}^p \eta_k^2 - \sum_{k=p+1}^{p+q} \eta_k^2,
\] 
где $p$ -- ПИИ, $q$ -- ОИИ(индексы инерции); $r=p+q$ -- ранг.

\subsection{Классификация ДУЧП2ПЛОСП в точке}
\begin{enumerate}[label=\Roman*.]
  \item \textbf{Эллиптический тип}
    \begin{itemize}[nosep]
      \item $r = n$ --- ранг $K$
      \item в каноническом виде все слагаемые одного знака 
            ($p = r = n,\ q = 0$ или $q = r = n,\ p = 0$)
      \item \textbf{\underline{Пример:}} уравнение Лапласа.
    \end{itemize}

  \item \textbf{Гиперболический тип}
    \begin{itemize}[nosep]
      \item $r = n$ --- ранг $K$
      \item в каноническом виде слагаемые имеют разные знаки 
            ($p \in [1, n-1]$ и $q \in [1, n-1]$)
      \item \textbf{Замечание:} при $p = n-1,\ q = 1$ или $p = 1,\ q = n-1$ --- 
            \emph{нормальный гиперболический} тип.
      \item \textbf{\underline{Пример:}} волновое уравнение.
    \end{itemize}

  \item \textbf{Параболический тип}
    \begin{itemize}[nosep]
      \item $r < n$ --- ранг $K$
      \item \textbf{Замечание:} при $p < n,\ q = 0$ или $p = 0,\ q < n$ --- 
            \emph{нормальный параболический} тип.
      \item \textbf{\underline{Пример:}} уравнение теплопроводности.
    \end{itemize}
\end{enumerate}

\begin{note}
Ранг $r$, ПИИ --- $p$, ОИИ --- $q$ не инвариантны для фиксированной точки $x^{(0)}$ (или функции фиксированной точки), поэтому тип ДУЧП2ПЛОСП зависит от выбора $x^{(0)} \in G$.
\end{note}

\textbf{\underline{Пример:}} уравнение Трикоми

\[
y \frac{\partial^2 u}{\partial x^2} + \frac{\partial^2 u}{\partial y^2} = 0
\]
Оно является \textbf{эллиптическим} при $y > 0$, \textbf{гиперболическим} при $y < 0$ и \textbf{параболическое} на линии вырождения $y = 0$.

\textbf{\underline{Пример:}} стационарное линеаризованное уравнение для потенциала скорости
\[
(1-M^2) \frac{\partial^2 u}{\partial x^2} + \frac{\partial^2 u}{\partial y^2} + \frac{\partial^2 u}{\partial z^2} = 0
\]
\begin{enumerate}[nosep]
  \item если \(M<1\) (дозвуковое течение), уравнение становится \textbf{эллиптическим};
  \item если \(M>1\) (сверхзвуковое течение), уравнение становится \textbf{гиперболическим};
  \item если \(M=1\) (звуковое течение), уравнение вырождается в \textbf{параболическое}.
\end{enumerate}

\begin{note}
Если в ДУЧП2ПЛОСП старшие коэффициенты $a_{ij} \neq a_{ij}(x)$(т.е. постоянные коэффициенты), то ЛНЗНП $\xi=\xi(x)$ можно привести это уравнение к уравнению канонического вида во всей $G$.
\end{note}

\begin{note}
Лапласиан в декартовой, сферической и цилиндрической системах координат.
\begin{enumerate}
% Прямоугольная
\item Декартовы координаты: $\displaystyle\Delta = \frac{\partial^2}{\partial x^2} + \frac{\partial^2}{\partial y^2} + \frac{\partial^2}{\partial z^2} 
$

% Цилиндрическая
\item Цилиндрические координаты: $\displaystyle\Delta = \frac{1}{r}\frac{\partial}{\partial r}\left(r \frac{\partial}{\partial r}\right) + \frac{1}{r^2}\frac{\partial^2}{\partial \varphi^2} + \frac{\partial^2}{\partial z^2}
$

% Сферическая
\item Сферические координаты: $\displaystyle\Delta = \frac{1}{\rho^2}\frac{\partial}{\partial \rho}\left(\rho^2 \frac{\partial}{\partial \rho}\right) + \frac{1}{\rho^2 \sin\theta}\frac{\partial}{\partial \theta}\left(\sin\theta \frac{\partial}{\partial \theta}\right) + \frac{1}{\rho^2 \sin^2\theta}\frac{\partial^2}{\partial \varphi^2}
$
\end{enumerate}
\end{note}
%FIXME НЕТ КАРТИИИИНОК
Попробуем обнулить коэффициенты $\tilde a_{kl}(\xi)$

\subsection{Характеристические поверхности(линии). Характеристики.}
\begin{definition}[Характеристические поверхности(линии). Характеристики.]
Пусть функция $\omega(x_1, \ldots, x_n) \in C^1(G)$ такова, что на плоскости $\omega(x)=0: \operatorname{grad}(\omega)  \neq0$.
Если на данной плоскости $\omega(x)$ удовлетворяет уравнению 
$\sum_{i, j=1}^n a_{ij}(x)\frac{\partial\omega}{\partial x_i}\frac{\partial\omega}{\partial x_j}=0$, то эта плоскость ($\omega(x)=0$) называется характеристической плоскостью (линией) или характеристикой для ДУЧП2ПЛОСП(1).
\end{definition}

Рассмотрим семейство характеристик $\omega(x, y)=C$, где $\alpha < C < \beta \Rightarrow$ некоторая достаточно малая область $G_\varepsilon$ будет целиком заполнена этими характеристиками, которые не будут пересекаться друг с другом, т.е. через каждую точку из $G_\varepsilon$ будет проходить единственная характеристика.

Если в НЗНП положить $\xi_1=\omega(x)$, то в $(\widetilde{1})$ коэффициент $\tilde a_{11} \equiv 0$.

Рассмотрим возможность приведения ДУЧП2ПЛОСП к каноническому виду в некоторой достаточно малой области $G_\varepsilon$. Для этого необходимо выполнение в $G_\varepsilon$ следующих условий:

\begin{quote}
$\tilde a_{kl}(\xi)=0,\quad k \neq l$

$\tilde a_{kk}(\xi)=\varepsilon_k\tilde a_{11}(\xi),\quad k=\overline{2, n},\ \varepsilon_k=0,\pm1$
\end{quote}

Всего наших соотношений:

\[
N=\frac{n^2-n}{2}+n-1 \leq n,
\]
где $n$ --- число независимых переменных.

Т.е. $n^2-n-2 \leq 0 \Rightarrow (n+1)(n-2) \leq0$, т.е. $n\leq2$

\textbf{\underline{Вывод:}} только лишь при $n \leq 2$ ДУЧП2ПЛОСП может быть приведено к каноническому виду в достаточно малой области с помощью НЗНП.

Рассмотрим случай $\boldsymbol{n=2}$.

\subsection{ДУЧП2ПЛОСП с двумя независимыми переменными}

Итак, \[a(x, y)\frac{\partial^2u}{\partial x^2} + 2b(x, y)\frac{\partial^2u}{\partial x \partial y} + c(x, y)\frac{\partial^2u}{\partial y^2} + \Phi(x, y, u, u_x, u_y)=0;\ (x, y) \in G \subset \mathbb{R}^2, a, b, c \in C(G)\ (1)\]

Рассмотрим НЗНП:
\begin{quote}
$\xi = \xi(x, y) \in C^2(G)$

$\eta = \eta(x, y) \in C^2(G)$
\end{quote}

Тогда в некоторой окрестности точки $\exists$ обратная замена \[x=x(\xi, \eta) \text{ и } y=y(\xi, \eta) \Rightarrow \tilde u(\xi, \eta)=u(x(\xi, \eta), y(\xi, \eta)) \text{ или } u(x, y)=u(\xi(x, y), \eta(x, y))\]

\noindent
Очевидно,
\begin{itemize}
  \item $\displaystyle 
    \frac{\partial u}{\partial x}
    = \frac{\partial \tilde u}{\partial \xi}\xi_x
    + \frac{\partial \tilde u}{\partial \eta}\eta_x
  $
  \item $\displaystyle 
    \frac{\partial u}{\partial y}
    = \frac{\partial \tilde u}{\partial \xi}\xi_y
    + \frac{\partial \tilde u}{\partial \eta}\eta_y
  $
  \item $\displaystyle 
    \frac{\partial^2 u}{\partial x^2}
    = \frac{\partial^2 \tilde u}{\partial \xi^2}\xi_x^2
    + 2\frac{\partial^2 \tilde u}{\partial \xi \partial \eta}\xi_x\eta_x
    + \frac{\partial^2 \tilde u}{\partial \eta^2}\eta_x^2
    + \frac{\partial \tilde u}{\partial \xi}\xi_{xx}
    + \frac{\partial \tilde u}{\partial \eta}\eta_{xx}
  $
  \item $\displaystyle 
    \frac{\partial^2 u}{\partial x \partial y}
    = \frac{\partial^2 \tilde u}{\partial \xi^2}\xi_x\xi_y
    + \frac{\partial^2 \tilde u}{\partial \xi \partial \eta}(\xi_x\eta_y + \xi_y\eta_x)
    + \frac{\partial^2 \tilde u}{\partial \eta^2}\eta_x\eta_y
    + \frac{\partial \tilde u}{\partial \xi}\xi_{xy}
    + \frac{\partial \tilde u}{\partial \eta}\eta_{xy}
  $
  \item $\displaystyle 
    \frac{\partial^2 u}{\partial y^2}
    = \frac{\partial^2 \tilde u}{\partial \xi^2}\xi_y^2
    + 2\frac{\partial^2 \tilde u}{\partial \xi \partial \eta}\xi_y\eta_y
    + \frac{\partial^2 \tilde u}{\partial \eta^2}\eta_y^2
    + \frac{\partial \tilde u}{\partial \xi}\xi_{yy}
    + \frac{\partial \tilde u}{\partial \eta}\eta_{yy}
  $
\end{itemize}


\section{Лекция 4: Продолжение}
Отсюда имеем:
\begin{equation}
	\tilde a \frac{\partial^2 \tilde u}{\partial \xi^2}+2\tilde b \frac{\partial^2 \tilde u}{\partial \xi\partial \eta}+
	\tilde c \frac{\partial^2 \tilde u}{\partial \eta^2} +
	\tilde{\Phi}(\xi,\eta,\tilde u, \frac{\partial\tilde u}{\partial \xi},\frac{\partial\tilde u}{\partial \eta}) = 0
	\tag{$\widetilde 1$}
\end{equation}
где:

\hspace*{2em} $\tilde a=a\xi^2_x+2b\xi_x\xi_y+c\xi^2_y$

\hspace*{2em} $\tilde b=a\xi_x\eta_x+b(\xi_x\eta_y+\xi_y\eta_x)+c\xi_y\eta_y$

\hspace*{2em} $\tilde c=a\eta^2_x+2b\eta_x\eta_y+c\eta^2_y$\\
Задача: выбрать такие $\xi$ и $\eta$, чтобы $\tilde a=\tilde c=0$, то есть:
$$ a\xi^2_x+2b\xi_x\xi_y+c\xi^2_y=0, \qquad (a\eta^2_x+2b\eta_x\eta_y+c\eta^2_y=0) $$
%ДИПСИЧНОЕ begin
\begin{note}
	Замена $(\xi,\eta)\leftrightarrow(x,y)$ предполагается невырожденной, т.е.
	\[
		J=\frac{\partial(\xi,\eta)}{\partial(x,y)} = \xi_x\eta_y - \xi_y\eta_x \neq 0.
	\]
	Это гарантирует существование обратного преобразования и эквивалентность уравнений.
\end{note}
%ДИПСИЧНОЕ begin
\begin{definition}[Характеристика и характерестическая кривая(случай $\mathbb{R}^2$, т.е. n=2)]
	Пусть функция $\omega(x,y) \in C^1(G):$ что на кривой $\omega(x,y)=0$ выполняется $\operatorname{grad}(\omega(x,y)) \neq \vec{0}$ и\\
	\begin{equation}
		a\omega^2_x+2b\omega_x\omega_y+c\omega^2_y=0
		\tag{2}
	\end{equation}
	Тогда кривая $\omega(x,y)=0$ называется характерестической кривой(характеристикой уравнения $(1)$), а само уравнение $(2)$ -- характерестическим уравнением уравнения $(1)$.
\end{definition}
\begin{lemma}[Связь первого интегралла и $\boldsymbol{\omega}$]%FIXME НАЗВАНИЕ КАКА
	Пусть $\omega(x,y) \in C^1(G)$ такова, что $\frac{\partial\omega}{\partial y}\ne 0$.\\
	Если $\omega(x,y)$ удовлетворяет уравнению $(2)$, то соотношение $\omega(x,y)=C$ является общим(первым) интеграллом ОДУ вида:
	\begin{equation}
		a\cdot\mathrm{d}y^2-2b\cdot\mathrm{d}x\mathrm{d}y+c\cdot\mathrm{d}x^2=0
		\tag{3}
	\end{equation}
	или $ a\cdot\mathrm{d}y^2-2b\cdot\mathrm{d}x\mathrm{d}y+c\cdot\mathrm{d}x^2=0 $, обозначим $(3')$\\
	Уравнение $(3)$ также называется характерестическим уравнением для $(1)$.
\end{lemma}
\begin{proof}
	Т.к. $\omega(x,y)\in C^1(G),\,\omega_y\ne 0$ и $\omega(x,y)$ удовлетворяет уравнению $(2)$, то из него получаем:\\
	$\displaystyle(2'):\; a\left(\frac{\omega_x}{\omega_y}\right)^2+2b\left(\frac{\omega_x}{\omega_y}\right)+c=0$ или
	$\displaystyle a\left(-\frac{\omega_x}{\omega_y}\right)^2-2b\left(-\frac{\omega_x}{\omega_y}\right)+c=0$
\end{proof}
Вспомним, что соотношение $\omega(x,y)=C$ является общим интегралом ОДУ, если это соотношение ($\omega(x,y)=C$) неявно задает функцию $y=y(x)$, которая
является решением указанного ОДУ.\\
Тогда по теореме ``О неявных функциях'', из $\omega_y\ne0\rightarrow$

\hspace*{1em}  1. Гиперболический тип: пусть в некоторой области $G$, $a\ge0$, $D>0$, тогда уравнения $(3)$ и $(3')$ -- различны и у них разлиные незвисимые 1\textsuperscript{ые} интеграллы: $\phi(x,y)=C_1$, $\psi(X,y)=C_2$ и $\displaystyle \frac{\partial(\phi,\psi)}{\partial(x,y)}\ne 0$\\%FIXME ЧТО ТУТ ПРОИСХОДИТ? ПЕРЕМЕШААААААЛИ??? короче структуру надо поправить на основе других лекций
Перейдем в уравнении $(1)$ к независимым переменным:
$\xi=\phi(x,y),\,\eta=\psi(x,y)$.\\
Тогда уравнение $(1)$ в некоторой окрестности из $G$ примет вид:

\hspace*{2em} $2\widetilde b \widetilde u_{\xi\eta}+\widetilde \Phi=0$, откуда $\widetilde u_{\xi\eta}=\widetilde{\widetilde \Phi}$ --- 1\textsuperscript{ая} каноническая форма.

\hspace*{1em}  2. Эллиптический тип: $\displaystyle \hat{u}_{\alpha\alpha}+\hat{u}_{\beta\beta}=\hat\Phi $;

\hspace*{1em}  3. Параболический тип: $\displaystyle \widetilde{u}_{\eta\eta}=\widetilde{\widetilde{\Phi}} $;
\begin{note}
	Если в некоторой области $G$ уравнение имеет гиперболический тип(ГТ), то в этой области $G$ существует 2 семейства $\mathbb R$(действительных) - характеристик, а если эллиптический тип, то в этой области $G$ существует 2 семейства $\mathbb C$(комплексно) - сопряженных характеристик.
\end{note}



\section{Лекция 5: Уравнения гиперболического типа(УГТ)}
Линейная замена: $\xi = \alpha + \beta,\ \eta = \alpha - \beta$ приводит исходное уравнение ко второй форме: $\hat{u}_{\alpha\alpha} - \hat{u}_{\beta\beta} = \hat{\Phi}$

\textbf{2. ЭТ}: пусть в некоторой области $G,\ a \ge 0,\ D < 0$ тогда $\phi(x, y)=C$ --- общий 1\textsuperscript{ый} интеграл  уравнения $(3_+^\prime) \Rightarrow \bar{\phi}(x, y) = \bar C$ --- тоже общий интеграл уравнения $(3_-^\prime)$.

Сделаем замену:

$\quad \xi=\phi(x, y),\ \eta = \bar\phi(x, y)$, получим: $\tilde u_{\xi \eta} = \tilde{\tilde{\Phi}}$.

Перейдем к $\mathbf{Re}$ - переменным:

$\quad \xi = \alpha + i\beta,\ \eta = \alpha - i\beta$, тогда имеем $\hat{u}_{\alpha\alpha} + \hat{u}_{\beta\beta} = \hat{\Phi}$

\textbf{3. ПТ}: пусть в некоторой области $G,\ a \ge 0,\ D = 0$. Пусть $b \ge 0$, тогда $b=\sqrt{a}\sqrt{c}$.

Уравнения $(3_+^\prime)$ и $(3_-^\prime)$ совпадают. Пусть $\phi(x, y)=C$ --- общий 1\textsuperscript{ый} интеграл уравнения $(3_+^\prime)$.

Пусть $\psi(x, y) \in C^2(G)$ --- произвольная функция, независимая с $\phi(x, y)$, т.е. $\frac{\partial(\phi,\ \psi)}{\partial(x,\ y)} \neq 0$.

Совершим НЗНП: $\xi = \phi(x, y),\ \eta=\psi(x, y) \overset{G}{\Rightarrow} \tilde a = 0 = (\sqrt{a}\xi_x+\sqrt{c}\xi_y)^2=0$, при этом \\ $\tilde b=(\sqrt{a}\xi_x + \sqrt{c}\xi_y)(\sqrt{a}\eta_x + \sqrt{c}\eta_y)=0\ \Rightarrow\ \tilde c \neq 0$ и уравнение $(1)$ примет вид: $\tilde c \tilde u_{\eta\eta} + \tilde \Phi = 0,$ или\\ $\tilde u_{\eta\eta} = \tilde{\tilde{\Phi}}$ --- 1\textsuperscript{ая} каноническая форма.

Аналогично, если $b \le 0$.

\subsection{Уравнения гиперболического типа(УГТ)}

Простейший случай:

$\quad u_{xx}-u_{yy}=0$ --- уравнение малых поперечны колебаний струны.

Процесс колебаний (малых поперечных) будем описывать поперечным смещением струны $u(x, t)$.

\begin{definition}[Струна]
Струной будем называть упругую гибкую нить(не сопротивляющуюся изгибу).
\end{definition}

Тогда натяжение струны в каждой точке направлено по касательной к мгновенному профилю струны в этой точке.

Малые поперечные колебания означают, что $|u| << l,\ l$ --- длина струны, $u_x^2 << 1$.

Тогда $\alpha \approx sin(\alpha) \approx tg(\alpha) \approx u_x << 1,\ cos(\alpha)=\frac{1}{\sqrt{1+u_x^2}} \approx 1$

Рассмотрим участок струны $(x_1,x_2)$:

\[
S^` = \int_{x_1}^{x_2} \sqrt{1+u_x^2}\,dx \approx \int_{x_1}^{x_2}\,dx = x_2 - x_1 = S = const
\]

Следовательно, $T(x, t) = T(x)$, т.е. натяжение струны $T(x, t)$ не зависит от времени $t$.

На самом деле, $T(x, t)=T_0=const$.

Действительно, при равновесии на $(x_1,\ x_2)$ вдоль $Ox$:
\[
T(x_2)cos(\alpha(x_2, t))=T(x_1)cos(\alpha(x_1, t)) \Rightarrow T(x_2) \approx T(x_1)
\](т.к. $cos(\alpha(x_2, t)) \approx cos(\alpha(x_1, t)) \approx 1$).

Выведем уравнение малых поперечных колебаний струны(УМПКС):

Второй закон Ньютона на $Ox$:
\[
T_x(x_2) \approx T(x_2) = T(x_1) = T_x(x_1) \Rightarrow T(x) \equiv T_0 = const
\]

Второй закон Ньютона на $Ou$: 
\[
\rho(x_1, t) \Delta x \frac{\partial^2 u}{\partial t^2}(x_1, t) \approx T_0 \frac{\partial u(x_2, t)}{\partial x} - T_0 \frac{\partial u(x_1, t)}{\partial x} + F(x_1, t) \Delta x;
\]
где $F(x_1, t) \Delta x$ --- распределенная нагрузка

\[
\rho(x_1, t) \frac{\partial^2 u}{\partial t^2}(x_1, t)
\approx \frac{T_0 \left[\frac{\partial u(x_2, t)}{\partial x} - \frac{\partial u(x_1, x)}{\partial x}\right]}{\Delta x} + F(x_1, t);\ \Delta x \to 0\
\]

\begin{equation}
\rho(x, t) \frac{\partial^2 u}{\partial t^2}
= T_0 \frac{\partial^2 u}{\partial x^2} + F(x, t);\ x \in (0; l)\ t > 0
\tag{1.F}
\end{equation}

\begin{itemize}[label={}, leftmargin=1.5cm, nosep]
  \item $\rho(x, t)$ --- линейная плотность струны
  \item $F(x, t)$ --- плотность поперечной нагрузки
  \item $T_0$ --- натяжение струны
\end{itemize}

В дальнейшем будем считать, что $\rho(x, t) = \rho_0 = const$.

Поделим на $\rho_0$, получим:

\begin{equation}
\frac{\partial^2 u}{\partial t^2} = 
a^2 \frac{\partial^2 u}{\partial x^2} + f(x, t), \quad x \in (0, l),\ t > 0,
\tag{1.f}
\end{equation} где
$
a = \sqrt{\frac{T_0}{\rho_0}} - \text{скорость};\ f(x, t) = \frac{F(x, t)}{\rho_0}
$

\textbf{Замечание.} Если в т. $x=x_0$ к струне приложена сосредоточенная сила $\Phi_0(t)$[$\delta-$функция с плотностью --- $\Phi(t) \delta(x-x_0)$], то в этой точке получим:
\[
T_0(u_x(x_0+0,t)-u_x(x_0-0,t)) = -\Phi_0(t)
\] или
\[
u_x(x, t)|_{x_0-0}^{x_0+0} = -\frac{\Phi_0(t)}{T_0}
\text{ --- излом струны в т. $x_0$.}
\]
\subsection{Простейшая смешанная задача уравнения колебаний струны(УКС) длины l, закрепленной на концах}

Неизвестная функция --- u(x, t)

\begin{equation}
\frac{\partial^2 u}{\partial t^2} = 
a^2 \frac{\partial^2 u}{\partial x^2} + f(x, t), \quad x \in (0, l),\ t > 0,
\tag{1.f}
\end{equation}

\textbf{Начальные условия(НУ):}

\begin{equation}
\begin{split}
u|_{t=0} &= u(x, 0) = u_0(x);\quad x \in [0, l] \ \text{(начальный профиль струны)} \\
u_t|_{t=0} &= u_t(x, 0) = u_1(x);\quad x \in [0, l] \ \text{(начальная скорость струны)}
\end{split}
\tag{2}
\end{equation}

\textbf{Однородные граничные условия 1 типа(ОГУ1Т):}

\begin{equation}
\begin{split}
u|_{x=0} &= u(0, t) = 0;\quad t \ge 0 \ \\
u|_{x=l} &= u(l, t) = 0;\quad t \ge 0 \ 
\end{split}
\tag{3.1.0}
\end{equation}

\textbf{Условия сшивки:}

\begin{equation}
\begin{split}
u_0(0) &= 0;(\text{т.к. }\ u_0(0)=u(0, 0)=0)\\
u_0(l) &= 0;(\text{т.к. }\ u_0(l)=u(l, 0)=0)\\
u_1(0) &= 0;(\text{т.к. }\ u_1(0)=u_t(0, 0)=0)\\
u_1(l) &= 0;(\text{т.к. }\ u_1(l)=u_t(l, 0)=0)\\
\end{split}
\tag{4.1.0}
\end{equation}

Введем $G=(0, l) \times (0, +\infty)$ и $\overline{G}=[0, l] \times [0, +\infty)$.

Решение ищем в классе $u \in C^2(G) \cap C^1(\overline{G})$

На $u_0(x),\ u_1(x),\ f(x,t)$ накладываем следующие требования:
\begin{itemize}[nosep, leftmargin=1.5cm]
  \item $u_0(x) \in C^2([0,l])$,
  \item $u_1(x) \in C^1([0,l])$,
  \item $f(x,t) \in C^1(\overline{G})$.
\end{itemize}

\begin{theorem}[Существования и единственности]
Решение простейшей краевой задачи для уравнения колебаний струны длины l, закрепленной на концах, $(1.f),\ (2),\ (3.1.0),\ (4.1.0)$ --- существует и единственно в рассматриевых классах функций.
\end{theorem}



\section{Лекция 6: Основные типы граничных условий(краевых).}
\input{sem1/lect6}

\section{Лекция 7: Редукция ЗК для УКБС}
\input{sem1/lect7}

\section{Лекция 9: Задача Гурса}
\subsection{Задача Гурса}
\begin{theorem}[существования и единственности]
	Решение $u(x, y) \in C^2(\Pi) \cap C^1(\overline{\Pi})$ задачи Гурса:
	\[
	Lu = u_{xy} + au_x + b_y + cu = f(x, y),\ (x, y) \in \Pi = (0, l) \times (0, h) \quad (1.f)
	\]
	\[
	\textbf{Условия на характеристиках(УХ): } u|_{x=0} = \phi(y),\ 0 \le y 
	\le h
	u|_{y=0} = \psi(x),\ 0 \le x \le l 
	\]
	\[
	\textbf{Условия согласования(УС): } \phi(0) = \psi(o) = u(0, 0),
	\] 
	\[\text{если } a(x, y) \in C^1(\overline{\Pi}),\ b(x, y), c(x, y), f(x, y) \in C^1(\overline{\Pi}), \phi(y) \in C^2([0, h]),\ \psi \in C^2([0, l]),
	\] существует и единственно.
\end{theorem}

Физический смысл: "сопло Лаваля".

\subsection{Смешанная задача(СЗ) для однородного уравнения колебаний полубесконечной струны(ОУКПБС), закрепленной на левом конце.}

\begin{align*}
	& u_{tt} = a^2 u_{xx}, \quad x > 0,\ t > 0 \tag{1.0} \\
	\textbf{НУ: } \quad & u|_{t=0}=u_0(x),\ u_t|_{t=0}=u_1(x),\ x \ge 0 \tag{2} \\
	\textbf{ОГУ1Т: } \quad & u|_{x=0}=0,\ t \ge 0 \tag{3.1.0} \\
	\textbf{УС: } \quad & u_0(0)=u_1(0)=0=u_0^{''}(0), \tag{4.1.0}
\end{align*}

$u(x, t) \in C^2((0, +\infty) \times (0, \infty)) \cap C^1([0, +\infty) \times [0, \infty))$

$u_0(x) \in C^2(0, +\infty),\quad u_1(x) \in C^1(0, +\infty)$

Вспоминим Задачу Коши(ЗК) для ОУКБС:
\begin{align*}
	& u_{tt} = a^2 u_{xx}, \quad x > 0,\ t > 0 \tag{1.0} \\
	\textbf{НУ: } \quad & u|_{t=0}=u_0(x),\ u_t|_{t=0}=u_1(x),\ x \in \mathbb{R} \tag{2}
\end{align*}
$u(x, t) \in C^2(t>0) \cap C^1(t \ge 0),\quad u_0(x) \in C^2(\mathbb{R}),\quad u_1 \in C^1(\mathbb{R})$

\begin{lemma}
	Если НУ в ЗК (1.0), (2) являются нечётными функциями относительно т. $x=x_0$, то соответствующее решение $u(x, t)$ ЗК (1.0), (2) равно нулю в т. $x=x_0$, т.е. $u(x_0, t) = 0,\quad t \ge 0$.
\end{lemma}

\begin{lemma}
	Если НУ в ЗК (1.0), (2) являются чётными функциями относительно т. $x=x_0$, то частная производная соответствующего решение $u(x, t)$ ЗК (1.0), (2) равна нулю в т. $x=x_0$, т.е. $u_x(x_0, t) = 0,\quad t \ge 0$.
\end{lemma}

\begin{note}
	функция $y = f(x)$, определенная не $(a, b)$ и симметричная относительно точки $x=x_0 \in (a, b)$ называется чётной(нечётной), если $f(2x_0 - x) = f(x)(f(2x_0 - x) = -f(x))$, где $f(x_0 - x^{'}) = \pm f(x_0 + x^{'})$(относительно т.$x_0$).
\end{note}

\begin{proof}
	(Доказательство первой леммы)
	Положим $x_0 = 0$. Тогда по теореме о существовании и единственности решения ЗК (1.0), (2) для ОУКБС в рассматриваемых классах функций $\exists!$ и представляется Ф.Д.0(формулой Даламбера):
	\[
		u(x, t) = \frac{1}{2} [u_0(x-at) + u_0(x+at)] + \frac{1}{2a}\int_{x-at}^{x+at}u_1(\xi)d\xi \Rightarrow
	\]
	\[
		u(0, t) = \frac{u_0(-at)+u_0(at)}{2} + \frac{1}{2a} \int_{-at}^{at} u_1(\xi)d\xi = 0,
	\]
	при $t \ge 0$, т.к. из условий леммы функции $u_0(x)$ и $u_1(x)$ --- нечётные. 
	
	Аналогично, доказывается вторая лемма с учётом того факта, что если $f(x)$ --- чётная функция, то её производная $f^{'}(x)$ --- нечётная.
\end{proof}

\subsection{Метод продолжений}

\textbf{Следствие 1} Для того, чтобы свести СЗ (1.0), (2), (3.1.0), (4.1.0) для ОУКПБС, закрепленной на левом конце, к ЗК для ОУКБС нужно начальные функции $u_0(x)$ и $u_1(x)$ нечетным образом продолжить влево за т. $x=0$.

\textbf{Следствие 2} Для того, чтобы свести СЗ для ОУКПБС, со свободным левым концом, к ЗК для ОУКБС нужно начальные функции $u_0(x)$ и $u_1(x)$ четным образом продолжить влево за т. $x=0$.

\begin{theorem}
	Решение $u(x, t)$ СЗ (1.0), (2), (3.1.0), (4.1.0) для ОУКПБС, закрепленной на левом конце, существует и единственно и представляется обобщенной Ф.Д.0(Ф.Д.0.1):
	
\end{theorem}

\end{document}
